\sscp{Uneven distribution in the lexicon}{07-02-02}
The lexicon is also revealing,
with different words being shared by different links in the chain.
It should also be noted that not all words that have been reconstructed
for Proto-Central Maluku (PCM) have the same scope,
as they encompass various combinations of the first three subgroups
listed in \frf{07-08}, or nodes within them.\footnote{
		Various PCM reconstructions are
		found in \citet{Stresemann1927,Collins1981,Collins1983,GrimesGrimes2020}.
		They must all be re-examined,
		as some should be limited to single subgroups,
		or even branches within single subgroups,
		and do not span all nodes of \tsc{Sula-Buru}, \tsc{Ambon-Seram},
		and \tsc{Seram-Tanimbar-Bomberai}.
		Others appear to. Some go beyond these three subgroups.
		All of these reflect the inadequacies of basing subgrouping on
		the lexicon.}
Many purported PCM reconstructions are found only in \tsc{Ambon-Seram} languages.
There are many good reasons for not basing subgroups on lexical
data, particularly in a region of long-term intensive contact,
extensive trade, and internal migrations (see chapters
\ref{ch:LayAdd}, \ref{ch:OthSub}, \ref{ch:Mar}, \ref{ch:MorLes}.

\begin{figure}[h]\centering
	\includegraphics[width=\columnwidth]{CM-LexicalWave}
	\fcp{Combinations in the shared lexicons of central Maluku}{07-08}
\end{figure}

While we find some words that appear to be exclusive to each
subgroup, we also find some words that are found in the following
combinations shown in \frf{07-08}.
Here \tsc{East Seram}, and \tsc{Eastern Islands} are separated
from the \tsc{Tanimbar-Bomberai} of \tsc{Seram-Tanimbar-Bomberai}
to illustrate things happening in the lexicon.

In some cases words are found in many languages in both subgroups.
In other cases a word may be found in most languages in one subgroup,
but only one or two non-adjacent languages in another,
such as in example \qf{07-01}.

{\wg\st{6pt}\begin{xltabular}[l]{\linewidth}{>{\hspace{-\tabcolsep}}l<{\hspace{-\tabcolsep}}>{\hspace{-\tabcolsep}\enspace}lll>{\raggedright\arraybackslash}X}
%{\wg\st{6pt}\tblr{>{\hspace{-\tabcolsep}}l<{\hspace{-\tabcolsep}}>{\hspace{-\tabcolsep}\enspace}llll}
\lsptoprule
\tbx{07-01}&& Word	& gloss & note \\ \midrule\hhline{~}
	&					&	**muyu	&	\tsc{neg}, no, not	&	(post-predicate in all languages)	\\
	&	Mangole	&	{\hr\hr}moyu	&	\tsc{neg}, no, not	&	\tsc{Sula-Buru}	\\
	&	Sanana	&	{\hr\hr}moya	&	\tsc{neg}, no, not	&	\tsc{Sula-Buru} (**u > \it{a}, see \srf{09-01})	\\
	&	Lisela	&	{\hr\hr}moo	&	\tsc{neg}, no, not	&	\tsc{Sula-Buru}	\\
	&	Buru		&	{\hr\hr}moo	&	\tsc{neg}, no, not	&	\tsc{Sula-Buru} (\it{tewa moo} ‘not know)’	\\
	&	N. Alune	&	{\hr\hr}mo	&	\tsc{neg}, no, not	&	\tsc{Ambon-Seram} (\it{tekwa mo} ‘dumb, not know’)\\ %\hhline{~}
%	&						&							&											&		\\
\lspbottomrule
\end{xltabular}\addtocounter{table}{-1}}

In some cases etyma (or irregular sound changes) systematically
skip over adjacent subgroups in the chain,
as in example \qf{07-02}.

{\wg\st{6pt}\tblr{>{\hspace{-\tabcolsep}}l<{\hspace{-\tabcolsep}}>{\hspace{-\tabcolsep}\enspace}llll}
\lsptoprule
\tbx{07-02}	& Subgroup & Proto-form & Gloss & Note \\ \midrule
	&	\tsc{Sula-Buru}		&*waŋka		&boat, canoe&\\
	&\tsc{Ambon-Seram}	&*baŋkaq	&boat, canoe&\\
	&\tsc{East Seram}		&*waŋka		&boat, canoe& node of STB\\
\lspbottomrule
\end{tabular}\mbox{}\\}

Some words have even broader scope stretching from Buru to Kowiai,
as in example \qf{07-03}.

{\wg\st{6pt}\begin{xltabular}[l]{\linewidth}{>{\hspace{-\tabcolsep}}l<{\hspace{-\tabcolsep}}>{\hspace{-\tabcolsep}\enspace}lll>{\raggedright\arraybackslash}X}
\lsptoprule
\tbx{07-03}	&	Language					&	{\hr\hr}Word	&	Gloss	&	Note		\\	\midrule
\endfirsthead
\lsptoprule
\qf{07-03}	&	Language					&	{\hr\hr}Word	&	Gloss	&	Note		\\	\midrule
\endhead
	&					&	**ntibu	&	fly (v.) & \citet[55]{Collins1983}, PCM\\
	&	Lisela	&	{\hr\hr}tifo, tifu	&	fly (v.)	&	\tsc{Sula-Buru}	\\
	&	Ambelau	&	{\hr\hr}en-tiβo	&	\tsc{3sg}-fly (v.)	&	\tsc{Sula-Buru}	\\
	&	Kayeli	&	{\hr\hr}tibo, en-tifo	&	fly (v.)	&	\tsc{Ambon-Seram}	\\
	&	Larike	&	{\hr\hr}tihu	&	fly (v.)	&	\tsc{Ambon-Seram}	\\
	&	Asilulu	&	{\hr\hr}tihu, rihu	&	fly (v.)	&	\tsc{Ambon-Seram} (inflected variants)	\\
	&	Kaitetu	&	{\hr\hr}kihu	&	fly (v.)	&	\tsc{Ambon-Seram} (inflected variants; reg. *nt > \it{k})	\\
	&	S. Wemale	&	{\hr\hr}dipu	&	fly (v.)	&	\tsc{Ambon-Seram} (reg. *nt > \it{d)}	\\
	&	N. Alune	&	{\hr\hr}ribu	&	fly (v.)	&	\tsc{Ambon-Seram} (reg. *nt > \it{r)}	\\
	&	Yalahatan	&	{\hr\hr}rihu	&	fly (v.)	&	\tsc{Ambon-Seram}	\\
	&	Manusela	&	{\hr\hr}rihu	&	fly (v.)	&	\tsc{Ambon-Seram}	\\
	&	Banda-Elat	&	{\hr\hr}tifu, ndifu	&	fly (v.)	&	Banda (inflected variants)	\\
	&	Masiwang	&	{\hr\hr}rihi	&	fly (v.)	&	STB, \tsc{East Seram}, (reg. *u > \it{i} /{\gap}{\#})	\\
	&	Kowiai	&	{\hr\hr}na-tiɸ	&	fly (v.)	&	STB, \tsc{Eastern Islands}	\\
\lspbottomrule
\end{xltabular}\addtocounter{table}{-1}}

Some words show both metathesised
and unmetathesised variants interspersed within the same subgroup
and non-contiguously within the wider region,
as in example \qf{07-04}.

{\wg\st{6pt}\begin{xltabular}[l]{\linewidth}{>{\hspace{-\tabcolsep}}l<{\hspace{-\tabcolsep}}>{\hspace{-\tabcolsep}\enspace}lll>{\raggedright\arraybackslash}X}
\lsptoprule
\tbx{07-04}	&	Language					&	{\hr\hr}Word	&	Gloss	&	Note		\\	\midrule
\endfirsthead
\lsptoprule
						&	Language					&	{\hr\hr}Word	&	Gloss	&	Note		\\	\midrule
\endhead
	&	PMP			&	\hr*buRah	&	\multicolumn{2}{l}{\parbox[t]{0.57\linewidth}{\raggedright
													spray water from the mouth; spray a mixture of saliva
													and masticated medicinal herbs on an ailing body part in curing}}	\\
	&					&	**tubuRa	&	\multicolumn{2}{l}{spit, spray masticated medicinal herbs}	\\
	&					&	**batuRa	&	\multicolumn{2}{l}{spit, spray masticated medicinal herbs}	\\	\hhline{~}
	&					&						&	&	\hp{\,}\\
	&	Taliabu	&	{\hr\hr}tafura-tsʸ	&	spit medicines	&	\tsc{Taliabu} (irr. *R > \it{r}; expect *R > \it{h})	\\
	&	Kadai	&	{\hr\hr}tafuɣa	&	spit medicines	&	\tsc{Taliabu} (irr. *R > \it{r}; expect *R > \it{h})	\\ \midrule\hhline{~}
	&	Lisela	&	{\hr\hr}stuha,	&	spit, spit medicines	&	\hp{\,}	\\	\hhline{~}
	&					&	{\hr\hr}ptuha,	&												&	\multirow{-2}{=}{\raggedright \tsc{Sula-Buru}	(\it{s{\tl}f}, body functions)}	\\
	&	Buru	&	{\hr\hr}eftuha	&	spit, spit medicines	&	\tsc{Sula-Buru}	\\
	&	Ambelau	&	{\hr\hr}erbuha	&	spit	&	\tsc{Sula-Buru} (reg. *t > \it{r})	\\  \midrule\hhline{~}
	&	Kayeli	&	{\hr\hr}egbula	&	spit	&	\tsc{Ambon-Seram}	\\
	&	Leliali	&	{\hr\hr}na-bula	&	\tsc{3sg}-spit	&	\tsc{Ambon-Seram} (dialect of Kayeli)	\\
	&	Boano	&	{\hr\hr}tuhuna	&	spit	&	\tsc{Ambon-Seram}	\\ \hhline{~}
	&				&									&				&	(reg. *R > \it{n})	\\
	&	Batumerah	&	{\hr\hr}irahula	&	spit	&	\tsc{Ambon-Seram}	\\
	&	Luhu	&	{\hr\hr}tahula-ye	&	spit	&	\tsc{Ambon-Seram}	\\
	&	Larike	&	{\hr\hr}tuhuda	&	spit	&	\tsc{Ambon-Seram}	\\
	&	Allang	&	{\hr\hr}pa-tahula	&	\tsc{caus}-spit	&	\tsc{Ambon-Seram}	\\
	&	Paulohi	&	{\hr\hr}tuhula	&	spit	&	\tsc{Ambon-Seram}	\\
	&	Amahai	&	{\hr\hr}suhura	&	spit	&	\tsc{Ambon-Seram}	\\
	&	Kaibobo	&	{\hr\hr}hula	&	spit medicine	&	\tsc{Ambon-Seram}	\\
	&	Kamarian	&	{\hr\hr}hura	&	spit medicine	&	\tsc{Ambon-Seram}	\\
	&	Haruku	&	{\hr\hr}tahura	&	spit	&	\tsc{Ambon-Seram}	\\
	&	Sepa	&	{\hr\hr}pa-tahula	&	\tsc{caus}-spit	&	\tsc{Ambon-Seram}	\\
	&	Tamilou	&	{\hr\hr}pa-tahula	&	\tsc{caus}-spit	&	\tsc{Ambon-Seram}	\\
	&	Tehoru	&	{\hr\hr}pa-tahula	&	\tsc{caus}-spit	&	\tsc{Ambon-Seram}	\\
	&	S. Wemale	&	{\hr\hr}tuhula	&	spit	&	\tsc{Ambon-Seram}	\\
	&	N.Wemale	&	{\hr\hr}tupula	&	spit	&	\tsc{Ambon-Seram}	\\
	&	Yalahatan	&	{\hr\hr}tuhula	&	spit	&	\tsc{Ambon-Seram}	\\
	&	Atamano	&	{\hr\hr}tufula	&	spit	&	\tsc{Ambon-Seram}	\\
	&	Manusela	&	{\hr\hr}fitoi (v.),	&	spit,	&	\tsc{Ambon-Seram}	\\ \hhline{~}
	&						&	{\hr\hr}hito (n.)	&	saliva	&	\hp{ }	\\
	&	Liana-Seti	&	{\hr\hr}au-bito	&	spit	&	\tsc{Ambon-Seram}	\\ \midrule\hhline{~}
	&	Banda	&	{\hr\hr}mbifiru	&	spit	&	Banda	\\ \midrule\hhline{~}
	&	Masiwang	&	{\hr\hr}fetoof,&	spit	&	STB, \tsc{East Seram}	\\	\hhline{~}
	&			&	{\hr\hr}batof	&	\hp{ }	&	\\
	&	Kei	&	{\hr\hr}ftuhut	&	spit	&	STB, \tsc{Tanimbar- Bomberai}	\\
	&	Irarutu	&	{\hr\hr}təɸí	&	spit	&	STB, \tsc{Tanimbar- Bomberai}	\\
\lspbottomrule
\end{xltabular}\addtocounter{table}{-1}}